\documentclass[12pt]{article}
\usepackage[utf8]{inputenc}
\usepackage[T1]{fontenc}
\usepackage{hyperref}
\title{Tambora Language Resources for K-12 Curriculum Development}
\author{Compiled by Claude Sonnet 5.5 (claude-sonnet-5-5) for Andr\'as Kornai}
\date{}
\begin{document}
\maketitle

\section{Introduction}
Tambora (ISO 639-3 xxt) was spoken in central Sumbawa, Indonesia, and died out with the Tambora culture in the 1815 eruption of Mount Tambora. The only substantial record is one word list collected before the eruption and published by Raffles in 1817 and 1830. Its classification is uncertain. Wikipedia lists it as unclassified, and Hammarström treats it as an isolate. A 2007 \emph{Oceanic Linguistics} article argues that it was a Papuan language. No native speakers, grammar or teaching materials were found.

\section{Primary Educational Resources}
No primers, textbooks or school programmes for Tambora were found.

\section{Linguistic and Reference Materials}
\begin{itemize}
\item Wikipedia, Tambora language: \url{https://en.wikipedia.org/wiki/Tambora_language}
\item Article on Tambora as a Papuan language (\emph{Oceanic Linguistics}, Project MUSE): \url{https://muse.jhu.edu/article/228921}
\end{itemize}

\section{Religious and Community Resources}
No religious or community-language resources were found.

\section{Digital and Audio Resources}
No audio or digital language resources were found. Only a historical word list exists.

\section{Curriculum Development Context}
The record is too thin for language instruction. A school unit could use Tambora to teach the 1815 eruption, the lost maritime trading culture of Sumbawa, and the problem of reconstructing a language from one word list. Local Sumbawa schools could pair it with the living languages of the island. The searches turned up no revitalization effort and no descendant community programme.
\end{document}